\chapter{Wstęp}
%======================================================================================================
\section{Wprowadzenie}

Wraz z rozwojem technologii, na rynku pojawia się coraz więcej innowacyjnych wynalazków wykorzystujących dotychczasowe osiągnięcia techniki informacyjnej i przyczyniających się do dalszego postępu technologicznego. Bardzo wyraźnym przykładem wspomnianego zjawiska jest przemysł informatyczny. Z upływem lat coraz wyraźniej zaobserwować można obecność urządzeń elektronicznych i komputerowych w różnych dziedzinach codziennego życia, począwszy od edukacji, przez opiekę zdrowotną, aż po relaks i rozrywkę. Nieodzowną rolę odgrywa tu rozwój Internetu, pozwalający na połączenie urządzeń w ogromne sieci do przesyłu informacji. 

Wiele dostępnych obecnie produktów na rynku konsumenckim stanowi wypadkową szerzącej się automatyzacji i coraz większych korzyści płynących z mediów takich jak Internet. Dobrze obrazuje to dziedzina tzw. Internetu Rzeczy (ang. \emph{Internet of Things}, \emph{IoT}), złożonego z wielu rozproszonych modułów sprzętowych podłączonych do wspólnej sieci komputerowej w celu stworzenia większego systemu.

% dopisać przegląd literatury
%======================================================================================================
\section{Cel i zakres pracy}

Celem pracy jest dokonanie przeglądu metod realizacji interfejsu sieciowego w modułach IoT oraz praktyczna implementacja (sprzętowa i programowa) wybranej metody w systemie bazującym na mikrokontrolerze z procesorem Cortex-M.

Zakres pracy obejmował:

\begin{itemize}
    \item [$\bullet$] przegląd możliwych do realizacji metod implementacji interfejsu sieciowego w modułach IoT bazujących na mikrokontrolerach z procesorami Cortex-M z uwzględnieniem m.in. wad i zalet każdej z nich,
    \item [$\bullet$] praktyczna realizacja modułu IoT w oparciu o mikrokontroler z procesorem Cortex-M z interfejsem sieciowym zrealizowanym według jednej z wybranych metod.
    
\end{itemize}
%======================================================================================================
\section{Struktura pracy}

Rozdział pierwszy pracy zawiera wstęp oraz cel i zakres, w ramach którego zrealizowana została niniejsza praca dyplomowa, a także skróconą informację o poszczególnych sekcjach pracy.

Celem drugiego rozdziału jest zapoznanie Czytelnika z pojęciem Internetu Rzeczy (IoT) wraz z jego formalną definicją, pokazując zalety oraz wady tej dziedziny. Rozdział zwieńczony został uzasadnieniem interpretacji modułów IoT jako urządzeń sieciowych.

Trzeci z rozdziałów omawia stos TCP/IP oraz jego relację w stosunku do modelu ISO OSI. Jego treść przedstawia i krótko opisuje protokoły dostępne w poszczególnych warstwach modelu TCP/IP. Ponadto, przedstawia on także różnicę pomiędzy transmisją TCP a transmisją UDP, oraz omawia sposób podłączenia warstwy fizycznej.

Tematyką czwartego rozdziału jest praktyczna implementacja interfejsu Ethernet oraz interfejsu bezprzewodowego w paśmie 2,4 GHz dla platformy sprzętowej Nucleo-F429ZI. Rozdział przedstawia krok po kroku sposób konfiguracji urządzenia oraz przygotowany kod źródłowy do obsługi interfejsu.

Ostatni rozdział omawia rezultaty pracy oraz przedstawia wnioski wyciągnięte w trakcie jej realizacji.